La definición de un cuerpo de agua abarca una gran cantidad de posibles ejemplos, entre ellos podemos tener los típicos como rios, lagos, mares y oceanos, pero también cosas como canales, pozos o tiros, presas, entre muchos otros.

El agua a su vez puede ser superficial o subterranea, y los cuerpos de agua pueden ser naturales o creados por el hombre.

El agua es potable cuando esta puede ser consumida sin hacer daño a la salud. Las aguas que no pueden ser consumidas son aguas residuales y agua de mar (que puede llegar hasta el 30\% de salinidad), así como aguas contaminadas.

Dos cuerpos de agua importantes son las lagunas y esteros, que se diferencian en lo siguiente: Las algunas son lagos de menor tamaño que cuentan con entrada de agua dulce y salada, donde la dulce llega por medio de rios y arroyos y la salada por filtración subterranea y mareas; por otro lado, los esteros solo cuentan con entrada de agua salada de manera directa.
Ambos son importantespara etapas iniciales de la fauna, donde las crias necesitan baja salinidad, o sea, que sea salobre, no salada.

Toda sustancia no necesaria para vivir y que afecta a la vida se conoce como xenobiótico, por ejemplo, plásticos, plaguicidas y desechos humanos en general. Todas estas sustancias son contanimantes, y as uvez pueden ser tóxicos para el ecosistema, caso en el que se les conocería como toxicoambientales. Algunos contaminantes pueden acumularse en un organismo, proceso llamado bioacumulación, y a su vez este proceso puede extrapolarse a lo largo de una cadena trófica, donde el proceso pasa a ser llamado biomagnificación.

Como biotecnólogos, nuestra responsabilidad es detectar contaminantes, buscar alternativas a lo que origina el contaminante y remediad el daño causado por este. La realización de todos estos análisis está normado.

Otros conceptos importantes son los canales de llamada y de salida, siendo el primero la entrada de agua marina a un estanque, y el segundo  representando la salida de agua del estanque.

\section{Contaminantes del agua}

Un cuerpo de agua es toda acumulación de agua, ya sea que esté en movimiento o no, así como si es natural o de origen humano.

Las sustancias externas a un medio (xenobióticos) que además causen daño a este son llamados contaminantes.
Todo contaminante debe contar con límites permisibles establecidos, y son valores máximos que ese contaminante debe reportar en un lugar dado. Los límites relacionados a un contaminante pueden variar dependiendo del cuerpo de agua del que se hable.
En caso de patógenos, por ejemplo Vibrio, Salmonella, etcétera, el límite es 0, o sea, no deben de existir. Otro ejemplo de contaminantes es el de los metales pesados.

Un proceso importante a prevenir es la eutroficación, proceso donde una alta concentración de nutrientes, muchas veces proveniente de una gran cantidad de materia orgánica, favorece la proliferación de microorganismos que generan una reducción en la concentración de oxígeno en el medio. Esto puede estar relacionado a descargas de materia fecal.
Las cadenas tróficas son un aspecto muy importante en el manejo ambiental, y la eutroficación afecta a las bases de estas. Este proceso afecta también a la zona fótica, que es la zona en que llega la luz solar, y cambios en esta afecta a la capacidad de los productores primarios de realizar sus actividades.
También son posibles cambios en el pH, esto por la actividad microbiana.

Otro ejemplo de contaminante son los microplásticos, que pueden ser disrruptores hormonales que afectan a la vida marina. Estos pueden ser también usados intencionalmente.

Los parámetros que se pueden medir en un análisis de contaminantes en cuerpos de agua pueden ser físicos, tales como los Sólidos disueltos/suspendidos, Materia flotante, Color, Olor, pH, Conductividad eléctrica; o Químicos, como el Oxígeno disuelto, Demanda química deoxígeno (DQO), Demanda bioquímica de oxígeno (DBO), Metales pesados, Nitrógeno, Fósforo, Cloruros	y los Fenoles.

Algunos fenómenos que sirven para el análisis de contaminantes en cuerpos de agua son la coagulación, Flotación y la Floculación.